\documentclass[11pt]{article}

\usepackage[margin=1in]{geometry}
\usepackage{amsmath}
\usepackage{listings}
\usepackage{xcolor}
\usepackage{hyperref}
\usepackage{booktabs}

\lstset{
  basicstyle=\ttfamily\small,
  breaklines=true,
  frame=single,
  columns=fullflexible,
  backgroundcolor=\color{gray!10}
}

\title{Natural Language Processing: Notes}
\author{ML A-Z --- Part 7}
\date{}

\begin{document}
\maketitle

\section{What this notebook implements}

The notebook \texttt{natural\_language\_processing.ipynb} builds a
sentiment classifier for restaurant reviews using a classic
\textbf{Bag of Words} pipeline:

\begin{enumerate}
  \item \textbf{Load the data.} Read \texttt{Restaurant\_Reviews.tsv},
        a tab-separated file of 1000 reviews, each labeled \texttt{0}
        (negative) or \texttt{1} (positive).
  \item \textbf{Clean the text.} For each review:
        \begin{itemize}
          \item strip everything except letters (\texttt{re.sub('[\^{}a-zA-Z]', ' ', \dots)}),
          \item lowercase and tokenize by splitting on whitespace,
          \item remove English stopwords (\texttt{nltk.corpus.stopwords}),
                \emph{except} the word \texttt{"not"}, which is kept because
                negation carries sentiment information,
          \item stem each remaining token with the Porter stemmer
                (e.g. \texttt{"loved"} $\to$ \texttt{"love"}).
        \end{itemize}
  \item \textbf{Vectorize.} Turn the cleaned corpus into a Bag-of-Words
        matrix with \texttt{CountVectorizer(max\_features=1500)}: each
        review becomes a row, each of the 1500 most frequent remaining
        words becomes a column, and each entry is that word's count in
        that review.
  \item \textbf{Split.} 80/20 train/test split (\texttt{random\_state=0}).
  \item \textbf{Classify.} Train a \texttt{GaussianNB} (Gaussian Naive
        Bayes) classifier on the training matrix.
  \item \textbf{Evaluate.} Predict on the test set and report a
        confusion matrix and accuracy score.

\end{enumerate}

\subsection{Result}
On the held-out test set (200 reviews), the model achieves:
\[
  \text{Confusion matrix} =
  \begin{pmatrix} 55 & 42 \\ 12 & 91 \end{pmatrix},
  \qquad \text{Accuracy} = 0.73
\]
i.e.\ 55 true negatives, 91 true positives, 42 false positives, and
12 false negatives, for 73\% overall accuracy.

\section{About the delimiter}

The dataset is loaded with:
\begin{lstlisting}[language=Python]
dataset = pd.read_csv('Restaurant_Reviews.tsv',
                       delimiter = '\t', quoting = 3)
\end{lstlisting}

\subsection{\texttt{delimiter = '\textbackslash t'}}
\texttt{Restaurant\_Reviews.tsv} is a \textbf{TSV} (tab-separated
values) file, not a comma-separated CSV. Review text routinely
contains commas (e.g.\ ``Great food, terrible service''), so a
comma delimiter would incorrectly split a single review into multiple
fields. Using \texttt{'\textbackslash t'} (the tab character) as the
field separator avoids this, since tabs essentially never appear
inside the review text itself.

\subsection{\texttt{quoting = 3}}
\texttt{quoting = 3} corresponds to \texttt{csv.QUOTE\_NONE} in
Python's \texttt{csv} module. It tells the parser to treat the quote
character (\texttt{"}) as an ordinary character instead of a field
delimiter. This matters here because review text frequently contains
stray double quotes (e.g.\ someone quoting the menu, or writing
\texttt{"amazing"} sarcastically). With the default quoting behavior,
pandas would try to interpret those quotes as marking the start/end of
a quoted field, which can corrupt parsing (merging rows or truncating
fields). Disabling quote interpretation makes the tab character the
only thing that matters for splitting columns, which is exactly what a
raw TSV of free-text reviews needs.

The four valid values of the \texttt{quoting} parameter (from Python's
\texttt{csv} module) are:
\begin{center}
\begin{tabular}{@{}cll@{}}
\toprule
Value & Constant & Meaning \\
\midrule
0 & \texttt{QUOTE\_MINIMAL} & Quote only fields containing special characters (default) \\
1 & \texttt{QUOTE\_ALL} & Quote every field \\
2 & \texttt{QUOTE\_NONNUMERIC} & Quote all non-numeric fields \\
3 & \texttt{QUOTE\_NONE} & Never treat the quote character specially \\
\bottomrule
\end{tabular}
\end{center}

\end{document}
